\section{Fast, precision-safe Cauchy feature application}
\label{sec:transforms}

\subsection{A different local interface}
The global theorem above reduces the cost of summing mode indices. This
section improves a different step: evaluating all source linear forms at a
given port assignment. Earlier packets give quadratic feature-table
construction or quadratic streamed application \cite{joint,kernels,cachedtree}.
The output of a full $N\times N$ table already has quadratic size. A
matrix-vector product has only linear input and output size and need not
materialize that table. We exploit this distinction.

Fast polynomial products, multipoint evaluation, transposed algorithms,
and Toeplitz convolution are classical \cite{mca}. The contribution here
is their exact realization for the corpus's normalized two-adic features,
including the extra precision required by the normalization. We found no
such implementation or argument in the inspected packet text. This is a
corpus-relative finding, not a claim of priority for fast structured linear
algebra in general.

Let $L=2^a$, $N=L/2$, $R=2^i$ with $0\le i<a$, and define the
increasing orientation slots
\[
 s_j=(j\bmod R)+2R\floor{j/R}\qquad(0\le j<N).
\]
Let $T$ be a specified element of valuation $i$ in the unramified
coefficient ring, available at the increased precision required below.
The source has $T=t+2^{a+1}\zeta$ with a fixed canonical description, so
such evaluation is available. For $\epsilon\in\{1,-1\}$ put
$x_j=\epsilon s_j$, $z_b=\epsilon s_b$ and
\[
 E_b=2\floor{b/R}-\pc(\floor{b/R}).
\]
The normalized features of \cite{joint} are
\begin{equation}\label{eq:features}
 F_{jb}=2^{-E_b}\frac{T}{T+x_j-z_b}
             \prod_{h<b}\frac{x_j-z_h}{T+x_j-z_h}.
\end{equation}
They are integral and vanish for $j<b$. Positive sign gives the target
features; negative sign gives the source features, commonly denoted $G$.
Thus the source forms are precisely $G^{\mathsf T}v$ for the vector of
real column bits. Omitted ports may be omitted from the sum or zero-filled.
We allow the first $1\le B\le N$ modes and any specified set $J$ of real ports.
Vector entries may belong to the finite free nilpotent extension as well.
Canonical lifts of vector entries can be chosen freely: integrality of
$F$ makes the output modulo $2^p$ independent of those choices. This does
not assert that canonical lifting is a ring homomorphism.

\subsection{The denominator valuations}
Let
\[
 D_B(X)=\prod_{h<B}(T+X-z_h),\qquad H=E_{B-1}.
\]
The sequence $E_b$ is nondecreasing. For $j\in J$, the known valuation of
the denominator is
\begin{equation}\label{eq:denvaluation}
 \kappa_j=v_2(D_B(x_j))
   =\sum_{d=1}^{i}\#\{h<B:h\equiv j\pmod{2^d}\}.
\end{equation}
Indeed the slot difference has valuation either below $i$ or strictly
above $i$; both slots have bit $i$ equal to zero. Since $v_2(T)=i$,
the summand denominator has valuation $\min(i,v_2(j-h))$, with $i$ for
$j=h$. Counting divisibility levels proves \eqref{eq:denvaluation}.
For $\kappa=\max_{j\in J}\kappa_j$, with zero for an empty set,
\[
 \kappa\le B(1-1/R)+i,\qquad H\le2(B-1)/R.
\]
In particular,
\begin{equation}\label{eq:guards}
 H+\kappa\le2B+i,\qquad H+2\kappa\le2B+2i.
\end{equation}
The implementation uses the exact valuations, not merely these upper bounds.

\subsection{Synthesis by a rational Newton product tree}
For a consecutive interval $I=[l,r)$ define
\[
 A_I(X)=\prod_{h\in I}(X-z_h),\quad
 B_I(X)=\prod_{h\in I}(T+X-z_h),
\]
and, with $w_b=c_b2^{H-E_b}T$,
\[
 C_I(X)=\sum_{b\in I}w_b
    \prod_{l\le h<b}(X-z_h)
    \prod_{b<h<r}(T+X-z_h).
\]
For consecutive halves $I=I_1\cup I_2$ the identity
\begin{equation}\label{eq:producttree}
 C_I=C_{I_1}B_{I_2}+A_{I_1}C_{I_2}
\end{equation}
is division free. At the root,
\[
 (Fc)_j=\frac{C_{[0,B)}(x_j)}{2^H D_B(x_j)}.
\]
Build the product tree and evaluate the numerator and denominator at all
$x_j$ modulo $2^{p+H+\kappa}$. Monic product/remainder trees need only
division by monic polynomials, so no nonunit leading coefficient is
inverted. By integrality of $F$, the numerator at $x_j$ is divisible
by $2^{H+\kappa_j}$. Divide that power of two exactly, invert the unit
part of the denominator, and reduce modulo $2^p$. The guard leaves at
least $p$ correct bits. All unit inverses can be obtained by one batch
inverse and linear multiplication work.

\subsection{The transpose requires a separate guard}
It would be invalid simply to transpose an algorithm that contains
divisions by powers of two. We instead derive a division-free linear
map and justify its final divisions directly.

Put $P=p+H+\kappa$. First build and evaluate $D_B$ at precision
$P_{\rm high}=p+H+2\kappa$. Its normalized units
$u_j=D_B(x_j)/2^{\kappa_j}$ are then known modulo $2^P$.
Batch-invert them and set
\[
 w_j=2^{\kappa-\kappa_j}u_j^{-1}v_j\pmod{2^P}.
\]
For $N_b(X)=\prod_{0\le h<b}(X-z_h)\prod_{b<h<B}(T+X-z_h)$, compute
$q_b=\sum_jN_b(x_j)w_j$ as follows. First obtain the first $B$ moments
\[
 m_k=\sum_jw_jx_j^k
\]
from the first $B$ coefficients of $\sum_jw_j/(1-x_jZ)$. A balanced
numerator/denominator construction and power-series inversion do this.
The series denominator has constant coefficient one. Newton inversion
$g\leftarrow g(2-fg)$ works in this ring; it does not require that two
be invertible.

Next transpose the division-free recurrence \eqref{eq:producttree}.
If $d$ is the parent coefficient functional, its child functionals are
\[
 d_1[k]=\sum_t B_{I_2}[t]d[k+t],\qquad
 d_2[k]=\sum_t A_{I_1}[t]d[k+t].
\]
Each is a polynomial correlation. At the leaves the resulting values
are the $q_b$. Finally,
\begin{equation}\label{eq:transposefinal}
 (F^{\mathsf T}v)_b=\frac{Tq_b}{2^{\kappa+E_b}}\pmod{2^p}.
\end{equation}
The numerator is divisible by the displayed power because $F$ is
integral. Since $E_b\le H$, the core precision $P$ leaves at least $p$
bits. The higher precision for denominator preparation is essential:
normalizing a denominator first consumes $\kappa_j$ bits before its
unit inverse is used at precision $P$.

\begin{theorem}[Fast guarded feature application]\label{thm:transform}
Let $m=B+|J|$. Both synthesis and the transposed source-form map can be
computed without a dense feature table in
\[
 O\bigl(\mathsf M_{\mathcal A}(m,p+O(B+i))\log(m+2)\bigr)
\]
arithmetic cost, plus linear unit arithmetic and a batch inverse, where
$\mathsf M_{\mathcal A}$ is polynomial-multiplication cost over the actual
coefficient algebra. A stored product tree uses $O(m\log(m+2))$
coefficient storage. The required precisions are respectively
$p+H+\kappa$ and $p+H+2\kappa$ for synthesis and transpose preparation.
\end{theorem}
\begin{proof}
The preceding identities prove correctness. At each balanced-tree level
the total polynomial degree is $O(m)$; standard superlinear
multiplication-cost bounds sum to the displayed complexity. Monic
remainders, constant-one series inversion, and correlations have the
same multiplication-based bounds. The storage bound counts the
coefficients across tree levels. Equations~\eqref{eq:guards} give the
precision assertion.
\end{proof}

\subsection{Odd interpolation nodes need only one convolution}
When $i=0$ the formula simplifies substantially: $s_j=2j$ and $T$ is a
unit. Let
\[
 f_j=v_2(j!)=j-\pc(j),\quad o_j=j!/2^{f_j},\quad
 D_j=\prod_{h=0}^j(T+2\epsilon h),\quad D_{-1}=1,
\]
and put $u_j=o_j/D_j$, $v_j=D_{j-1}/o_j$. All these denominators are
units. Direct cancellation in \eqref{eq:features} gives
\begin{equation}\label{eq:toeplitz}
 F_{jb}=\epsilon^bT\,
  2^{f_j-f_b-f_{j-b}}u_jv_{j-b}\qquad(j\ge b).
\end{equation}
The exponent equals $v_2\binom jb\ge0$. To verify the identity, the
numerator product in \eqref{eq:features} is
$(2\epsilon)^b j!/(j-b)!$, the denominator product is
$D_j/D_{j-b-1}$, and $E_b=b+f_b$.

\begin{theorem}[One-convolution source forms]\label{thm:oneconv}
For all $N$ modes and ports at $i=0$, either $Fc$ or $F^{\mathsf T}v$
is computable with one polynomial convolution, $O(N)$ preparation,
and one batch inverse at precision $p+2f_{N-1}$.
\end{theorem}
\begin{proof}
Set $K=f_{N-1}$ and $V_t=2^{K-f_t}v_t$ at that precision.
For synthesis convolve $V$ with
$c'_b=c_b\epsilon^bT2^{K-f_b}$. Multiply coefficient $j$ by $u_j$ and
divide by $2^{2K-f_j}$; equation~\eqref{eq:toeplitz} gives $(Fc)_j$.
For the transpose correlate $V$ with $a_j=v^{\rm input}_j2^{f_j}u_j$.
At index $b$, multiply by $\epsilon^bT$ and divide by $2^{K+f_b}$.
Again \eqref{eq:toeplitz} gives exactly the desired answer. One reversed
polynomial product yields all correlations. Each division is exact and
consumes at most $2K$ bits. The odd parts $o_j$ and products $D_j$ are
prepared by recurrence and batch inversion in linear work.
\end{proof}

In the scalar ring, quasi-linear integer multiplication with Kronecker
substitution gives $\widetilde O(N(p+N))$ bit operations, including
the guard. It is therefore $\widetilde O(Np)$ when $p\ge N$. The
coefficient algebra's rank must be included for extension rings, or one
can leave the cost as $\mathsf M_{\mathcal A}$. The executable supplied
here implements the scalar ring with CPython integers. It does not
implement an optimal FFT backend or an optimized extension-ring version.

Taking $p=p_t$ at the active-node scale $p\asymp MD/u$ and $N=RD/2$ gives
$p/N\asymp 2M/(Ru)$. Under the fixed schedule of this report,
$R\le2^u=s^{\eta_d+o(1)}$, while active $M>s/\lambda^{3/4}$.
Thus the extra $O(N)$ bits are lower order than the active precision.
This argument must be rechecked for a different schedule; it is not a
license to omit guards from a general complexity claim.

Finally, these transforms evaluate linear source forms. The subsequent
nonlinear mode monomials and the paid child contractions still have to be
summed. The transforms alone therefore leave the global exponent scale
unchanged. Their benefit is a provable local asymptotic improvement and,
for the odd-node specialization, a measured improvement on the supplied
small kernel fixtures.
